% SPDX-License-Identifier: Apache-2.0
\section{A Reservation Station}
\label{sec:x-station}

A reservation station gathers the operands of an operation that
arrive out of order and on separate channels, and fires when all of
them have. \code{Station2} to \code{Station10} in \code{txhdl\_parts}
are the stations of two to ten inputs, each written out by
\code{station!}. An input carries a tag and a value, a
\code{Tagged<TB, T>}; the station has a line per tag value, one cell
per input in each, and an input lands in its cell of the line its
tag names. A line whose every cell holds is sent on the output as
one value, the tag and the values; at most one line per cycle, the
completing input of lowest index choosing when two complete at once,
the value that completes a line going straight through and the rest
coming from their cells. An input's \code{ready} is its take: its
cell is free, and taking it either completes nothing or completes
the line being sent this cycle. The tutorial on the
station~\cite{station} reads its rule cycle by cycle.

The example is a station of two inputs, a byte and a half word, by a
two-bit tag. The two sources offer their tags in different orders,
one with gaps, and the sink holds the output off two cycles in four,
so one line waits for room and a repeated tag waits for its cell.
The netlist is checked against this run under both simulators, and
the same station written by hand in Verilog, \code{station.v}, runs
beside it as a foreign unit fed the same offers, the two compared
line by line; the tutorial~\cite{station} reads the two side by
side.

\subsection*{Side by side with Verilog}

The rows below set the hand-written Verilog beside the TxHDL text
\code{station!} wrote for the same two-input station, piece by
piece: the ports, the state, what each input offers, what completes
a line, what is sent and taken, the drives, and the outputs. The two
are the same hardware; the Verilog names every width and slices the
tag and the value by hand, and writes each \code{ready} and
\code{valid}, where the TxHDL text takes the widths from the payload
type and the handshake from \code{recv\_if} and \code{send}. Both
are included from the tree, and the example runs the two on the same
offers and asserts they agree.

\newcommand{\sidebyside}[2]{%
\noindent\textbf{#2}\par\nopagebreak\vspace{2pt}%
\noindent\begin{minipage}[t]{0.49\linewidth}
\lstinputlisting[style=ascii,%
  basicstyle=\ttfamily\fontsize{5.6}{6.6}\selectfont,%
  rangebeginprefix=//\ begin\{,rangebeginsuffix=\},%
  rangeendprefix=//\ end\{,rangeendsuffix=\},%
  includerangemarker=false,linerange=#1-#1]{lib/examples/station.v}
\end{minipage}\hfill
\begin{minipage}[t]{0.49\linewidth}
\lstinputlisting[%
  basicstyle=\ttfamily\fontsize{5.6}{6.6}\selectfont,%
  rangebeginprefix=//\ begin\{,rangebeginsuffix=\},%
  rangeendprefix=//\ end\{,rangeendsuffix=\},%
  includerangemarker=false,linerange=#1-#1]{station2_marked.rs}
\end{minipage}\par\vspace{6pt}}

\sidebyside{ports}{The ports: Verilog, left; TxHDL, right.}
\sidebyside{state}{The state: the cells and the masks.}
\sidebyside{reads}{What each input offers, and whether its cell is free.}
\sidebyside{completes}{What completes a line.}
\sidebyside{takes}{The line sent, and what is taken.}
\sidebyside{drives}{The drives: the cells and the masks.}
\sidebyside{outputs}{The outputs: the line, and its handshake.}

\begin{figure}[htbp]
\centering
\input{station_timing.tex}
\caption{The station: the two inputs' tags, values and handshakes,
the occupancy masks, and each line as it goes out, its tag and its
two values.}
\label{fig:x-station}
\end{figure}

\lstinputlisting[caption={\code{ex\_station.rs}. Run with \code{bazel run //lib/examples:ex\_station}.},label={lst:x-station}]{lib/examples/ex_station.rs}

\lstinputlisting[style=ascii,caption={What \code{ex\_station} printed when the build ran it: what each input offered, the masks, and the line taken per cycle, then the lowering.},label={lst:x-station-out}]{out_station.txt}
